\chapter{OpenROAD Integration with eSim}

This chapter describes the complete implementation work carried out during the integration of OpenROAD with eSim. It explains the frontend modifications, backend automation modules, netlist conversion process, OpenROAD Flow Scripts execution, and the debugging steps performed during development.

\section{Architecture Overview}

The integration connects the eSim simulation environment with the OpenROAD RTL-to-GDSII backend flow. The overall workflow implemented during this internship is shown below:

\begin{enumerate}
    \item The user selects a project inside eSim and clicks the \textbf{OpenROAD} button.
    
    \item The button handler inside \texttt{Application.py} identifies the selected project and locates the generated \texttt{.cir.out} file.
    
    \item The \texttt{netlist2rtl.py} module parses the \texttt{.cir.out} file and generates a synthesizable Verilog (\texttt{.v}) file.
    
    \item The \texttt{OpenROAD.py} backend module automatically generates the required ORFS configuration files such as \texttt{config.mk} and \texttt{constraint.sdc}.
    
    \item OpenROAD Flow Scripts execute synthesis, floorplanning, placement, CTS, routing, and GDSII generation.
    
    \item The generated layout is visualized automatically using the OpenROAD GUI and KLayout.
\end{enumerate}

\section{Modification of Application.py}

The eSim frontend is primarily controlled through the file \texttt{src/frontEnd/Application.py}. During this internship, a new OpenROAD execution button was added to the left-side toolbar of the graphical interface.

The following code snippet shows the implementation used for adding the OpenROAD button:

\begin{lstlisting}[language=python,
caption={Adding the OpenROAD button inside Application.py.}]
self.openroadAction = QtWidgets.QAction(
    QtGui.QIcon(os.path.join(
        self.resource_path,
        'openroad.png')),
    'OpenROAD',
    self
)

self.openroadAction.setStatusTip(
    'Run OpenROAD Flow'
)

self.openroadAction.triggered.connect(
    self.openOpenROAD
)

self.leftToolBar.addAction(
    self.openroadAction
)
\end{lstlisting}

The button was connected with a backend execution handler that validates the selected project and launches the OpenROAD flow.

\section{OpenROAD Button Handler}

The handler function retrieves the currently selected project directory, validates the generated \texttt{.cir.out} file, and initializes the backend execution module.

\begin{lstlisting}[language=python,
caption={OpenROAD execution handler in Application.py.}]
def openOpenROAD(self):

    project_dir = self.obj_projectExplorer.curProject

    if not project_dir:
        QtWidgets.QMessageBox.warning(
            self,
            'No Project',
            'Please select a project.'
        )
        return

    project_name = os.path.basename(project_dir)

    cir_out = os.path.join(
        project_dir,
        project_name + '.cir.out'
    )

    if not os.path.isfile(cir_out):
        QtWidgets.QMessageBox.warning(
            self,
            'File Not Found',
            'Run KiCad-to-Ngspice conversion first.'
        )
        return

    from maker.OpenROAD import OpenROADFlow

    self.openroad_obj = OpenROADFlow(
        cir_out,
        project_dir
    )

    self.openroad_obj.run()
\end{lstlisting}

Figure~\ref{fig:button_gui} shows the OpenROAD button integrated into the eSim GUI toolbar.

\begin{figure}[H]
\centering
\includegraphics[width=0.50\textwidth]{OpenROAD-Button.png}
\caption{OpenROAD button integrated into the eSim GUI toolbar.}
\label{fig:button_gui}
\end{figure}

\section{Implementation of netlist2rtl.py}

The \texttt{netlist2rtl.py} module was developed to automatically convert eSim-generated \texttt{.cir.out} files into synthesizable Verilog RTL.

The parser reads the SPICE netlist line by line and extracts:

\begin{itemize}
    \item module name,
    \item input ports,
    \item output ports,
    \item subcircuit instances,
    \item internal net connections.
\end{itemize}

\subsection{Netlist Parsing}

The following logic was used for parsing the \texttt{.cir.out} file:

\begin{lstlisting}[language=python,
caption={Core parsing logic in netlist2rtl.py.}]
def parse_cir_out(cir_out_path, project_dir):

    module_name = os.path.basename(project_dir)

    ports_in = []
    ports_out = []
    instances = []

    with open(cir_out_path, 'r') as f:

        for line in f:

            line = line.strip()

            if line.startswith('*') or not line:
                continue

            if line.lower().startswith('v'):

                parts = line.split()

                if len(parts) >= 3:

                    net = parts[1]

                    if net.lower() not in ('gnd', '0'):
                        ports_in.append(net)

            if line.lower().startswith('x'):
                instances.append(line)

    return (
        module_name,
        ports_in,
        ports_out,
        instances
    )
\end{lstlisting}

\subsection{Verilog Generation}

After parsing, the extracted information is converted into a Verilog module.

\begin{lstlisting}[language=python,
caption={Verilog generation in netlist2rtl.py.}]
def generate_verilog(
    module_name,
    ports_in,
    ports_out,
    instances,
    output_path
):

    lines = []

    all_ports = ports_in + ports_out

    lines.append(
        f'module {module_name} ('
    )

    lines.append(
        '    ' + ',\n    '.join(all_ports) + ');'
    )

    for p in ports_in:
        lines.append(f'input {p};')

    for p in ports_out:
        lines.append(f'output {p};')

    for inst in instances:
        lines.append(' ' + inst + ';')

    lines.append('endmodule')

    with open(output_path, 'w') as f:
        f.write('\n'.join(lines) + '\n')
\end{lstlisting}

\section{Implementation of OpenROAD.py}

The \texttt{OpenROAD.py} module serves as the primary backend controller for ORFS execution.

The module performs the following tasks:

\begin{itemize}
    \item validates input files,
    \item generates ORFS configuration files,
    \item creates timing constraints,
    \item launches ORFS execution,
    \item verifies generated outputs,
    \item opens the generated layout in OpenROAD GUI.
\end{itemize}

\subsection{config.mk Generation}

The following configuration file was generated automatically for each project:

\begin{lstlisting}[language=python,
caption={config.mk generation inside OpenROAD.py.}]
def generate_config_mk(
    design_name,
    verilog_path,
    output_dir
):

    config = f"""
export DESIGN_NICKNAME = {design_name}
export DESIGN_NAME     = {design_name}
export PLATFORM        = sky130hd

export VERILOG_FILES   = {verilog_path}

export SDC_FILE = \
./designs/$(PLATFORM)/\
$(DESIGN_NICKNAME)/constraint.sdc

export CORE_UTILIZATION  = 40
export CORE_ASPECT_RATIO = 1
export CORE_MARGIN       = 2
export PLACE_DENSITY     = 0.65
"""

    config_path = os.path.join(
        output_dir,
        'config.mk'
    )

    with open(config_path, 'w') as f:
        f.write(config.strip())

    return config_path
\end{lstlisting}

\subsection{SDC File Generation}

Timing constraints were automatically generated using the following logic:

\begin{lstlisting}[language=python,
caption={Automatic SDC generation.}]
def generate_sdc(
    design_name,
    ports_in,
    ports_out,
    output_path
):

    lines = [
        f'current_design {design_name}',
        'create_clock -name clk -period 10.0000'
    ]

    for p in ports_in:
        lines.append(
            f'set_input_delay 0.1000 '
            f'-add_delay [get_ports {{{p}}}]'
        )

    for p in ports_out:
        lines.append(
            f'set_output_delay 0.1000 '
            f'-add_delay [get_ports {{{p}}}]'
        )

    with open(output_path, 'w') as f:
        f.write('\n'.join(lines) + '\n')
\end{lstlisting}

\section{ORFS Execution}

Once configuration files were generated successfully, OpenROAD Flow Scripts were executed using the Makefile-based backend flow.

\begin{lstlisting}[language=python,
caption={ORFS execution from OpenROAD.py.}]
def run_orfs(design_name, orfs_path):

    cmd = [
        'make',
        f'DESIGN_CONFIG=./designs/sky130hd/'
        f'{design_name}/config.mk'
    ]

    result = subprocess.run(
        cmd,
        cwd=os.path.join(orfs_path, 'flow'),
        text=True
    )

    if result.returncode != 0:
        raise RuntimeError(
            f'ORFS failed for {design_name}'
        )
\end{lstlisting}

\section{Validation of .cir.out Files}

Before launching the OpenROAD flow, validation checks were added to verify that the generated \texttt{.cir.out} file existed and contained valid SPICE data.

\begin{lstlisting}[language=python,
caption={Validation of .cir.out files.}]
def validate_cir_out(cir_out_path):

    if not os.path.isfile(cir_out_path):
        return False, 'File not found.'

    with open(cir_out_path, 'r') as f:
        content = f.read().lower()

    if '.end' not in content:
        return False, 'Invalid .cir.out file.'

    return True, 'OK'
\end{lstlisting}

\section{Debugging and Fixes}

Several issues were encountered during implementation and testing. These issues were identified and resolved during the integration process.

\subsection{ORFS Path Errors}

Incorrect project naming initially caused ORFS configuration paths to fail. The issue was resolved by sanitizing project names before directory creation.

\begin{lstlisting}[language=python,
caption={Project name sanitization.}]
import re

design_name = re.sub(
    r'[^A-Za-z0-9_]',
    '',
    project_name.replace(' ', '_')
)
\end{lstlisting}

\subsection{HalfAdder Verilog Generation Issue}

The generated Verilog file for the HalfAdder project was initially empty because of case-sensitive module naming differences.

The issue was resolved by extracting the module name directly from the project directory instead of parsing it from the netlist.

\begin{lstlisting}[language=python,
caption={Module name extraction fix.}]
module_name = os.path.basename(project_dir)
\end{lstlisting}

\subsection{DEF Loading Error in OpenROAD GUI}

Initially, the OpenROAD GUI failed to load the generated DEF file correctly because LEF technology files were not loaded.

The issue was resolved by launching the GUI using the generated \texttt{.odb} database file.

\begin{lstlisting}[language=python,
caption={Corrected OpenROAD GUI launch.}]
def launch_openroad_gui(
    orfs_path,
    design_name
):

    result_dir = os.path.join(
        orfs_path,
        'flow',
        'results',
        'sky130hd',
        design_name
    )

    odb_file = os.path.join(
        result_dir,
        '6_final.odb'
    )

    cmd = [
        os.path.join(
            orfs_path,
            'tools',
            'install',
            'OpenROAD',
            'bin',
            'openroad'
        ),
        '-gui',
        odb_file
    ]

    subprocess.Popen(cmd)
\end{lstlisting}

\section{FullAdder and HalfAdder Validation}

The complete integration flow was tested using both \textbf{FullAdder} and \textbf{HalfAdder} circuits.

The integration successfully generated:

\begin{itemize}
    \item synthesized Verilog files,
    \item DEF floorplan files,
    \item GDSII layout files,
    \item timing constraint files,
    \item parasitic extraction outputs.
\end{itemize}

Figure~\ref{fig:both_projects} shows both projects visible inside the eSim project panel.

\begin{figure}[H]
\centering
\includegraphics[width=0.50\textwidth]{both.png}
\caption{FullAdder and HalfAdder projects inside the eSim workspace.}
\label{fig:both_projects}
\end{figure}

Figure~\ref{fig:fa_process} shows the FullAdder project during OpenROAD flow execution.

\begin{figure}[H]
\centering
\includegraphics[width=0.50\textwidth]{FA-Project-by-process.png}
\caption{FullAdder project during OpenROAD execution.}
\label{fig:fa_process}
\end{figure}